% =====================================================================
%  From trained transformers to B-RASP (annealing + rightmost UHAT).
%  Uses macros from transformers.tex.
% =====================================================================

\providecommand{\Vals}{\mathcal{V}}            % finite set of hidden-state values
\providecommand{\Sc}{S}                        % score table
\providecommand{\Outtab}{\delta}               % output table (transduction)

\section{From Trained Transformers to \texorpdfstring{$\brasp$}{B-RASP}}
\label{sec:trained_to_brasp}

\Cref{sec:transformers} translated $(1,1)$-precision transformers. Here we describe a route from
an ordinary trained transformer to a $\brasp$ autoregressor, in three steps.
First, we anneal the softmax temperature of attention, and optionally of the readout, until attention is
effectively hard (\cref{sec:annealing}).
Second, at inference time we replace the remaining softmax attention by unique-hard attention with rightmost tie-breaking (\cref{sec:rightmost}).
Third, we observe that each layer of the resulting model is a transduction between finite sets of
values, and we translate it into $\brasp$ layer by layer (\cref{sec:finite_transductions}).

\subsection{Annealing the temperature}
\label{sec:annealing}

We start from a standard transformer with real-valued weights, no position embeddings, and
future masking, as in \cref{def:transformer} but without rounding.
We add an inverse temperature $\beta>0$ to the attention and to the readout:
\begin{align*}
\alpha_{ij} &= \frac{\exp(\beta s_{ij})}{\sum_{j'\leq i}\exp(\beta s_{ij'})}
\quad\text{where } s_{ij} = \Wq(\mathbf{h}^{(\ell-1)}_i)\cdot\Wk(\mathbf{h}^{(\ell-1)}_j), \\
\symprob{\autoregressor_\beta}{\stru}{\tau} &= \softmax\bigl(\beta\,\Wout(\mathbf{h}^{(L)}_{|\stru|})\bigr)(\tau).
\end{align*}
The usual transformer is the case $\beta=1$.
As $\beta\to\infty$, that is, as the temperature $1/\beta$ goes to $0$,
$\softmax(\beta\mathbf{z})$ converges to $\avghardmax(\mathbf{z})$, the uniform distribution over $\argmax\mathbf{z}$.
So in the limit, every attention head averages the values at its maximal-score positions,
and the readout is the $\avghardmax$ distribution of \cref{eq:hardmax_symprob}.

Annealing the readout is optional. The translation accepts any readout of the final hidden state
(\cref{sec:hardmax_autoregressive}), so we may instead anneal only the attention and keep a softmax, quantized
softmax, or $\mathrm{top}_k$ readout. Annealing the readout to $\avghardmax$ gives the fewest distinct
probabilities, and hence the cheapest model counting, but it makes generation nearly greedy:
the readout is random only where logits tie.

To reach this limit without destroying the trained solution, we increase $\beta$ gradually
and keep training at each step. Starting from the trained weights at $\beta=1$,
we multiply $\beta$ by a constant factor at each stage, fine-tune at the new $\beta$ until the loss
recovers, and stop when the attention weights and the next-symbol distributions are within a
small tolerance of their hard limits on held-out data.
At that point, the model is close to its zero-temperature limit on the inputs we care about.

\subsection{Rightmost tie-breaking at inference}
\label{sec:rightmost}

The zero-temperature limit of attention is uniform over all positions attaining the maximal
score. Without position embeddings, ties are common: two positions holding the same hidden state
always receive the same score. Averaging over unboundedly many tied positions lets the model
count (for example, the fraction of $a$'s in the prefix), and $\brasp$ cannot express this.
So at inference we break ties by choosing the rightmost maximal position, as in
\citet{yang-etal-2024-masked}. This gives the following model.

\begin{definition} \label{def:rightmost_uhat}
A \defn{rightmost unique-hard attention transformer} has the same parameters as a transformer
(\cref{def:transformer}), over $\R$ instead of $\F$. On input $\str=w_1\cdots w_n$,
it computes $\mathbf{h}^{(0)}_i=E(w_i)$ and, for $\ell\in[L]$,
\begin{align}
j_i &= \max\Bigl\{j\leq i \Bigm| s_{ij}=\max_{j'\leq i}s_{ij'}\Bigr\}, &
\mathbf{h}^{(\ell)}_i &= f^{(\ell)}\bigl(\mathbf{h}^{(\ell-1)}_i+\Wv^{(\ell)}(\mathbf{h}^{(\ell-1)}_{j_i})\bigr),
\label{eq:rightmost}
\end{align}
where $s_{ij}=\Wq^{(\ell)}(\mathbf{h}^{(\ell-1)}_i)\cdot\Wk^{(\ell)}(\mathbf{h}^{(\ell-1)}_j)$.
Its next-symbol distribution is $\rho(\Wout(\mathbf{h}^{(L)}_{|\stru|}))$ for a readout $\rho$ from \cref{tab:readouts}, by default $\avghardmax$.
\end{definition}

The readout is unchanged, since $\brasp$ autoregressors can express any readout of the final hidden state.
The annealed model and the rightmost model agree at every position where the maximal attention
score is attained only once. Where it is attained several times, they may differ.
Whether this matters for a given model is an empirical question, which we check by
comparing the two models' next-symbol distributions on held-out data.

\subsection{Layers as finite transductions}
\label{sec:finite_transductions}

A rightmost unique-hard attention transformer can be simulated in $\brasp$ because each of its hidden states takes only finitely many values,
and these can be computed from the trained weights.
The embedding layer produces at most $|\alphabet|$ values. If layer $\ell-1$ produces values in a finite set
$\Vals_{\ell-1}$, then by \cref{eq:rightmost} the output of layer $\ell$ at $i$ depends only on the
values at $i$ and at the attended position $j_i$. So it lies in
\[
\Vals_0 := E(\alphabet), \qquad
\Vals_\ell := \bigl\{\, f^{(\ell)}\bigl(\mathbf{u}+\Wv^{(\ell)}(\mathbf{u}')\bigr) \bigm| \mathbf{u},\mathbf{u}'\in\Vals_{\ell-1} \,\bigr\}.
\]
Layer $\ell$ is therefore a length-preserving transduction from strings over $\Vals_{\ell-1}$ to
strings over $\Vals_\ell$, given by two finite tables:
\[
\Sc_\ell(\mathbf{u},\mathbf{u}') := \Wq^{(\ell)}(\mathbf{u})\cdot\Wk^{(\ell)}(\mathbf{u}'),
\qquad
\Outtab_\ell(\mathbf{u},\mathbf{u}') := f^{(\ell)}\bigl(\mathbf{u}+\Wv^{(\ell)}(\mathbf{u}')\bigr),
\qquad \mathbf{u},\mathbf{u}'\in\Vals_{\ell-1}.
\]
The score table $\Sc_\ell$ says how much a query holding $\mathbf{u}$ prefers a key holding
$\mathbf{u}'$. The output table $\Outtab_\ell$ gives the new state when $\mathbf{u}$ attends to $\mathbf{u}'$.
We compute both tables from the trained weights, in the same floating-point arithmetic used at
inference, so that ties in $\Sc_\ell$ are exactly the ties the model sees.

The sets can grow as $|\Vals_\ell|\leq|\Vals_{\ell-1}|^2$. In practice we keep only the values that
some input actually produces. Computing exactly which values occur is not needed:
any superset of the occurring values works, because formulas for values that never occur are never satisfied.
Multi-head layers are handled the same way, with one attended value $\mathbf{u}'$ per head.

\subsection{Translation to \texorpdfstring{$\brasp$}{B-RASP}}
\label{sec:uhat_to_brasp}

We now translate each layer, given as the tables $\Sc_\ell$ and $\Outtab_\ell$, into $\brasp$.
As in \cref{thm:transformer_to_brasp}, we build formulas $\phi_{\mathbf{h}^{(\ell)}=\mathbf{v}}$ that hold
exactly where the hidden state equals $\mathbf{v}$.
The attended position depends on the query value, so, as before, we enumerate the query value
$\mathbf{q}$. We also enumerate the maximal score $\gamma$ that the query attains.
Once $\mathbf{q}$ and $\gamma$ are fixed, the positions that the query may attend to are given by a unary formula,
and rightmost $\brasp$ attention finds the attended position directly.

Fix $\mathbf{q}\in\Vals_{\ell-1}$ and a score $\gamma$ in the row $\Sc_\ell(\mathbf{q},\cdot)$. By \cref{eq:pointwise},
the following unary formulas say that a position scores exactly $\gamma$, or more than $\gamma$, against $\mathbf{q}$:
\[
\phi^{\mathbf{q}}_{=\gamma} := \bigvee_{\mathbf{u}\colon\Sc_\ell(\mathbf{q},\mathbf{u})=\gamma}\phi_{\mathbf{h}^{(\ell-1)}=\mathbf{u}},
\qquad
\phi^{\mathbf{q}}_{>\gamma} := \bigvee_{\mathbf{u}\colon\Sc_\ell(\mathbf{q},\mathbf{u})>\gamma}\phi_{\mathbf{h}^{(\ell-1)}=\mathbf{u}}.
\]
At a position $i$ holding $\mathbf{q}$, the maximal score over $j\leq i$ is $\gamma$ iff some
$j\leq i$ scores $\gamma$ and none scores more:
\[
\phi^{\mathbf{q}}_{\max=\gamma} := \cntge{1}{\phi^{\mathbf{q}}_{=\gamma}(y)} \land \lnot\cntge{1}{\phi^{\mathbf{q}}_{>\gamma}(y)}.
\]
The attended position $j_i$ is then the rightmost $j\leq i$ satisfying $\phi^{\mathbf{q}}_{=\gamma}$. This is
$i$ itself if $i$ satisfies it, and otherwise the position found by strict rightmost attention.
So the attended value is $\mathbf{u}'$ iff
\[
\phi^{\mathbf{q},\gamma}_{\mathrm{att}=\mathbf{u}'} :=
\bigl(\phi^{\mathbf{q}}_{=\gamma}(x)\land\phi_{\mathbf{h}^{(\ell-1)}=\mathbf{u}'}(x)\bigr)\lor
\bigl(\lnot\phi^{\mathbf{q}}_{=\gamma}(x)\land
\attrdefault{y}{y<x}{\phi^{\mathbf{q}}_{=\gamma}(y)}{\phi_{\mathbf{h}^{(\ell-1)}=\mathbf{u}'}(y)}{\bot}\bigr).
\]

\begin{theorem} \label{thm:uhat_to_brasp}
Let $\tfr$ be a rightmost unique-hard attention transformer of depth $L$, with value sets
$\Vals_\ell$ and tables $\Sc_\ell,\Outtab_\ell$ as above. For every $\ell\leq L$ and
$\mathbf{v}\in\Vals_\ell$, the $\brasp$ formula $\phi_{\mathbf{h}^{(\ell)}=\mathbf{v}}$ defined by
\begin{align*}
\phi_{\mathbf{h}^{(0)}=\mathbf{v}} &:= \bigvee_{\sigma\colon E(\sigma)=\mathbf{v}}Q_\sigma, \\
\phi_{\mathbf{h}^{(\ell)}=\mathbf{v}} &:= \bigvee_{\substack{\mathbf{q},\mathbf{u}'\in\Vals_{\ell-1}\colon\\ \Outtab_\ell(\mathbf{q},\mathbf{u}')=\mathbf{v}}}
\phi_{\mathbf{h}^{(\ell-1)}=\mathbf{q}}\land\phi^{\mathbf{q}}_{\max=\Sc_\ell(\mathbf{q},\mathbf{u}')}\land\phi^{\mathbf{q},\Sc_\ell(\mathbf{q},\mathbf{u}')}_{\mathrm{att}=\mathbf{u}'}
\qquad(\ell\geq1)
\end{align*}
satisfies
\[
\str,i\models\phi_{\mathbf{h}^{(\ell)}=\mathbf{v}}\iff\mathbf{h}^{(\ell)}_i(\str)=\mathbf{v}.
\]
These formulas have total size $O\bigl(|\alphabet|\,|\Vals_0| + \sum_{\ell}|\Vals_{\ell-1}|^3\bigr)$.
\end{theorem}

\begin{proof}
By induction on $\ell$. The case $\ell=0$ is immediate. For $\ell\geq1$, let
$\mathbf{q}=\mathbf{h}^{(\ell-1)}_i$, let $\gamma$ be the maximal score of $\mathbf{q}$ over $j\leq i$, and let
$\mathbf{u}'=\mathbf{h}^{(\ell-1)}_{j_i}$.
By the induction hypothesis, $\phi^{\mathbf{q}}_{\max=\gamma'}$ holds at $i$ exactly for $\gamma'=\gamma$, and
$\phi^{\mathbf{q},\gamma}_{\mathrm{att}=\mathbf{u}''}$ holds at $i$ exactly for $\mathbf{u}''=\mathbf{u}'$.
Hence the only disjunct whose three conjuncts can all hold at $i$ is the one for $(\mathbf{q},\mathbf{u}')$.
So $\phi_{\mathbf{h}^{(\ell)}=\mathbf{v}}$ holds at $i$ iff $\Outtab_\ell(\mathbf{q},\mathbf{u}')=\mathbf{v}$, that is, iff $\mathbf{h}^{(\ell)}_i=\mathbf{v}$.
For the size, each layer has $|\Vals_{\ell-1}|^2$ disjuncts. Each disjunct refers to formulas
$\phi^{\mathbf{q}}_{=\gamma}$, $\phi^{\mathbf{q}}_{>\gamma}$ of size $O(|\Vals_{\ell-1}|)$, which are shared across disjuncts.
\end{proof}

Each layer uses one level of attention for the counting quantifiers and one for the rightmost
attention, so the $\brasp$ program has depth $O(L)$.
Finally, the readout is translated exactly as in \cref{cor:hardmax_to_brasp_ar}: for $\avghardmax$ and $\mathrm{top}_k$,
with one formula per output symbol and the uniform $\arout$; for other readouts, with one formula per distribution in $R$.
The result is a $\brasp$ autoregressor with the same next-symbol distributions as the rightmost
unique-hard attention transformer, and it can be passed directly to the encoding of \cref{sec:bounded-brasp}.